

%% 
%% This is a LaTeX document generated by automatic conversion of a
%% Mathematica notebook using Mathematica 4.0.
%% 
%% This document uses special macros that are defined in a LaTeX 
%% package file with a name of the form notebook*.sty.  Appropriate 
%% commands for loading the package have been included in this document.
%% The LaTeX package files can be found in a directory with path:
%% 
%% /usr/local/mathematica/SystemFiles/IncludeFiles/TeX/
%% 
%% The LaTeX package used in this document may require that your 
%% LaTeX implementation support fonts other than the standard Computer 
%% Modern fonts that are used by LaTeX.  See the Additional Information 
%% section under the TeXSave entry in the Mathematica Reference Guide 
%% for more information.
%% 
%% Implementation-specific instructions for installing TeX-related files
%% can be found at this URL.
%% 
%% http://www.wolfram.com/support/FrontEnds/Export/TeX
%% 
%% 

\documentclass{article}
\usepackage{notebook2e, latexsym}


\begin{document}

\Title{DISCRETE-TIME FOURIER TRANSFORM}
\Subtitle{LECTURE 15}
\Subsubtitle{ELE314: LINEAR SIGNALS AND SYSTEMS}
\Subsubtitle{Mariusz Jankowski \copyright{}

University of Southern Maine

mjkcc@usm.maine.edu}
\dispSFinmath{<<\Mvariable{SignalPlot`}}
\dispSFinmath{<<\Mvariable{Graphics`Graphics`}}
\begin{CellGroup}
\Section*{1 Introduction}


The discrete-time Fourier transform is a fundamental analysis tool in  discrete-time signals and systems theory. An understanding of the basic  concepts
of discrete-time Fourier synthesis and analysis is indispensable to  a communications or control professional, since the rapid advancement of  digital
computer technology and integrated circuit fabrication, has made  digital signal processing preferrable to analog, in numerous application  areas.

{\itshape "While continuous-time Fourier analysis can be traced to the  investigation of problems in mathematical physics in the 1700's, the tools
 for analyzing discrete-time signals and systems have their own distinct  roots. In particular, discrete-time concepts and methods are fundamental
to  the discipline of numerical analysis. Formulas for the processing of discrete  sets of data points to produce numerical approximations for interpolation,
 integration, and differentiation were being investigated as early as the time  of Newton in the 1600's... In the 1940s and 1950s major steps forward
were  taken in the development of discrete-time techniques in general and in the  use of the tools of discrete-time Fourier analysis in particular.
The reasons  for this push forward were the increasing use and capabilities of digital  computers and the development of design methods for sampled-data
systems,  that is, discrete-time systems for the processing of sampled continuous-time  signals." }[Oppenheim \&{} Willsky, 1983].

The discrete-time Fourier transform plays  the same role as its classic continuous-time counterpart. Conceptualy, the  transforms are same, both
decompose aperiodic signals into a continuum of  frequencies. However, the differences are striking and interesting. 




\end{CellGroup}
\begin{CellGroup}
\Section*{2 DTFT Pair of Equations}


The discrete-time Fourier transform is defined by the following pair of equations:

\inlineTFinmath{x[n]\multsp =\multsp \frac{1}{2\pi }\multsp \int _{-\pi }^{\pi }X(\theta )\multsp \multsp {{\ExponentialE }^{\multsp j\multsp \multsp
\theta \multsp \multsp n}}\DifferentialD \omega }



\inlineTFinmath{X(\theta )\multsp =\multsp \sum _{n=-\infty }^{\infty }x[n]\multsp \multsp {{\ExponentialE }^{-j\multsp \theta \multsp \multsp n}}}



where, \(\omega \)  is the radial frequency. The function \inlineTFinmath{X(\omega )} is referred to as the {\itshape discrete-time Fourier transform}
(DTFT) of \inlineTFinmath{x[n]}. (1) is known as the {\itshape inverse} transformation, while (2) is the {\itshape forward} transformation.



The first important difference between the  continuous-time and discrete-time Fourier transforms can be seen in the  analysis formula (2): it is
not an integral, but a summation. Another major  difference between the continuous and discrete cases is the periodicity of \inlineTFinmath{X(\theta
)}. \inlineTFinmath{X(\theta )} is periodic with period T\(=\)2\(\pi \), namely \inlineTFinmath{X(\theta )\multsp =\multsp X(\theta +2\pi )}. Recall
that the discrete-time complex exponential is a periodic function  in frequency, with a period of 2\(\pi \) (see \buttonbox[Hyperlink]{Lecture 2}).
The interval \inlineTFinmath{-\pi \leq \multsp \theta \multsp \leq \pi } is known as the fundamental interval.



\begin{CellGroup}
\Exercise{Problem 1}
\ExerciseText{Prove that \inlineTFinmath{X(\theta )\multsp =\multsp X(\theta +2\pi )}.}



\end{CellGroup}
We now consider the DTFT of the signal \inlineTFinmath{x[n]\multsp =\multsp {r^n}\multsp u[n]}{\itshape .} By direct evaluation using (2) we immediately
get



\begin{CellGroup}
\dispSFinmath{\multsp \sum _{n=0}^{\infty }{r^n}\exp [-\imag \multsp \theta \multsp n]}
\dispSFoutmath{\frac{{{\ExponentialE }^{\ImaginaryI \multsp \theta }}}{{{\ExponentialE }^{\ImaginaryI \multsp \theta }}-r}}

\end{CellGroup}
\dispSFinmath{X[\Mvariable{\theta \_}]:=\frac{{e^{\imag \multsp \theta }}}{-r+{e^{\imag \multsp \theta }}}}
For a value of \inlineTFinmath{r=0.5}, the magnitude and phase are as follows,



\begin{CellGroup}
\dispSFinmath{\MathBegin{MathArray}{l}
\Muserfunction{DisplayTogetherArray}\big[\big\{\Mfunction{Plot}\big[\Mfunction{Abs}\big[X[\theta ]/.\big\{r\rightarrow \frac{1}{2}\big\}\big],\{\theta
,-2\multsp \pi ,4\multsp \pi \}\big],  \\
\noalign{\vspace{1.17708ex}}
\hspace{2.em} \Mfunction{Plot}\big[\Mfunction{Arg}\big[X[\theta ]/.\big\{r\rightarrow \frac{1}{2}\big\}\big],\{\theta ,-2\multsp \pi ,4\multsp \pi
\}\big]\big\}\big];\\
\MathEnd{MathArray}}
\mathGraphic{lecture15.tex_gr1.eps}

\end{CellGroup}
The frequency response \inlineTFinmath{X(\theta )} of the discrete-time signal \inlineTFinmath{x[n]} is a complex and continuous function of frequency
\(\theta \).  In addition,  as stated above, the DTFT is periodic with period 2\(\pi \).  The synthesis  equation is therefore similar in form to
the analysis equation in  continuous-time Fourier series theory. This leads to the interesting notion  of duality between the CTFS and DTFT (see
Oppenheim/Willsky, pp. 395-396). 



We now turn our attention to the inverse transformation. We calculate the  unit sample response \inlineTFinmath{h[n]} given a filter with the following
frequency response:



\inlineTFinmath{H(\theta )\multsp =\multsp \big\{\MathBegin{MathArray}[c]{cc}
	1,\multsp &-\frac{\pi }{3}+2\pi \multsp l\leq \theta \leq \frac{\pi }{3}+2\pi \multsp l\multsp  \\
	0,&\Mvariable{else} 
  \MathEnd{MathArray}}

where l is an integer. Here is a plot of the frequency spectrum of the filter.



\mathGraphic{lecture15.tex_gr2.eps}
\begin{CellGroup}
Using (1) we get

\dispSFinmath{\frac{1}{2\pi }\int _{-\frac{\pi }{3}}^{\frac{\pi }{3}}\exp [\imag \multsp \theta \multsp n]\DifferentialD \theta \multsp //\Mfunction{ExpToTrig}}
\dispSFoutmath{\frac{\sin \big[\frac{n\multsp \pi }{3}\big]}{n\multsp \pi }}
which is the familiar sinc function.

\dispSFinmath{\MathBegin{MathArray}{l}
x[0]=\frac{1}{3};  \\
\noalign{\vspace{1.34375ex}}  \\
x[\Mvariable{n\_}]:=\frac{\sin \big[\frac{n\multsp \pi }{3}\big]}{n\multsp \pi };
\MathEnd{MathArray}}
\begin{CellGroup}
\dispSFinmath{\Muserfunction{LinePlot}[\Mfunction{Table}[\{n,x[n]\},\{n,-10,30\}],\Mvariable{PlotRange}->\Mvariable{All}];}
\mathGraphic{lecture15.tex_gr3.eps}

\end{CellGroup}

\end{CellGroup}
\begin{CellGroup}
\Exercise{Problem 2}
\ExerciseText{Obtain the inverse DTFT of two lowpass filters

(a) \inlineTFinmath{{G_1}(\theta )\multsp =\multsp H\big(\frac{\theta }{2}\big)}.

(b) \inlineTFinmath{{G_2}(\theta )\multsp =\multsp H(2\theta )}.

Compare all three results. How does the width of the filter affect the time-domain characteristics of the unit sample response of the filter?}



\end{CellGroup}

\end{CellGroup}
\begin{CellGroup}
\Section*{3 Frequency Response of LTI Discrete-Time Systems}


Earlier in the semester we considered the unit sample response of the  so-called moving average filter (MA). Recall that the impulse response of
an MA filter is a unit pulse scaled by its length N. It was verified experimentally that this filter acts as a smoothing operator. We will now consider
the frequency response of an N-th order MA filter and show that it is a low-pass filter. Hence the smoothing effect on data. Consider the rectangular
pulse of length \inlineTFinmath{N=2M+1} (odd length). 



\dispSFinmath{h[\Mvariable{n\_},\Mvariable{M\_}]:=\frac{1}{2\multsp M\multsp +1}\multsp \Mfunction{If}[\Mfunction{Abs}[n]\leq M,1,0];}
\begin{CellGroup}
\dispSFinmath{\MathBegin{MathArray}{l}
\Muserfunction{LinePlot}[\Mfunction{Table}[\{n,h[n,5]\},\{n,-10,20\}],  \\
\noalign{\vspace{0.5ex}}
\hspace{1.em} \Mvariable{PlotLabel}\rightarrow \Mvariable{Moving\multsp average\multsp filter}];\\
\MathEnd{MathArray}}
\mathGraphic{lecture15.tex_gr4.eps}

\end{CellGroup}
\begin{CellGroup}
The DTFT is given by the following function (try this calculation by hand):

\begin{CellGroup}
\dispSFinmath{\MathBegin{MathArray}{l}
H[\Mvariable{\theta \_}]:=\multsp \sum _{n=-5}^{5}\exp [-\imag \multsp \theta \multsp n]\multsp //\Mfunction{\multsp }\Mvariable{ExpToTrig};  \\
\noalign{\vspace{1.01042ex}}  \\
H[\theta ]
\MathEnd{MathArray}}
\dispSFoutmath{1+2\multsp \cos [\theta ]+2\multsp \cos [2\multsp \theta ]+2\multsp \cos [3\multsp \theta ]+2\multsp \cos [4\multsp \theta ]+2\multsp
\cos [5\multsp \theta ]}

\end{CellGroup}
\begin{CellGroup}
\dispSFinmath{\MathBegin{MathArray}{l}
\Mfunction{Plot}[\Mfunction{Abs}[H[\theta ]],\{\theta ,-\pi ,3\multsp \pi \},  \\
\noalign{\vspace{0.5ex}}
\hspace{1.em} \Mvariable{PlotRange}->\Mvariable{All},\Mvariable{PlotLabel}\rightarrow \Mvariable{DTFT\multsp of\multsp MA\multsp filter}];\\
\MathEnd{MathArray}}
\mathGraphic{lecture15.tex_gr5.eps}

\end{CellGroup}
Clearly, the filter passes low frequencies and attenuates high frequencies  in the signal. Two features of particular importance are the mainlobe
width and the stopband ripple. The mainlobe width is given by the location of the first zero, while the stopband ripple is defined by the maximum
deviation of the  frequency response from zero, measured at the peak of the first sidelobe. The  mainlobe width may be found by using FindRoot.




\end{CellGroup}
\begin{CellGroup}
\dispSFinmath{\Mfunction{FindRoot}[H[\theta ]==0,\{\theta ,0.3\}]}
\dispSFoutmath{\{\theta \rightarrow 0.571199\InvisibleSpace -5.68646\times {{10}^{-18}}\multsp \ImaginaryI \}}

\end{CellGroup}
while the peak sidelobe ripple will be estimated by plotting the normalized frequency response in the vicinity of the peak value.

\begin{CellGroup}
\dispSFinmath{\MathBegin{MathArray}{l}
\Mfunction{Plot}[\Mfunction{Abs}[H[\theta ]/H[0]],\{\theta ,0.5,1.\},\Mvariable{PlotRange}\rightarrow \{0,0.25\},  \\
\noalign{\vspace{0.5ex}}
\hspace{1.em} \Mvariable{GridLines}\rightarrow \{\Mvariable{None},\{\Mvariable{Automatic},\Mfunction{Hue}[2/3]\}\}];\\
\MathEnd{MathArray}}
\mathGraphic{lecture15.tex_gr6.eps}

\end{CellGroup}
\begin{CellGroup}
\Exercise{Problem 3}
\ExerciseText{(a) Plot the spectrum for \inlineTFinmath{M=10} and \inlineTFinmath{M=21}.

(b) Discuss how the spectrum of the finite, discrete-time pulse changes with increasing pulse width? Specifically, consider the mainlobe width and
sidelobe ripple.}



\end{CellGroup}

\end{CellGroup}
\begin{CellGroup}
\Section*{PROBLEM SOLUTIONS}
\begin{CellGroup}
\Exercise{Problem 1}
\ExerciseText{Prove that \inlineTFinmath{X(\theta )} is periodic with period T\(=\)2 \(\pi \), i.e. show that \inlineTFinmath{X(\theta )\multsp =\multsp
X(\theta +2\pi )}.}


\ExerciseText{{\bfseries Solution.}}


\ExerciseText{\inlineTFinmath{\MathBegin{MathArray}{l}
X(\theta +2\pi )=\multsp \sum _{n=-\infty }^{\infty }x[n]\multsp \multsp {{\ExponentialE }^{-j(\multsp \theta +2\pi )\multsp \multsp n}}=\multsp
\multsp   \\
\noalign{\vspace{1.40625ex}}
\hspace{2.em} {{\ExponentialE }^{-j\multsp 2\pi \multsp \multsp n}}\multsp \sum _{n=-\infty }^{\infty }x[n]\multsp \multsp {{\ExponentialE }^{-j\multsp
\theta \multsp n}}\multsp =\multsp X(\theta )\multsp \\
\MathEnd{MathArray}}}



\end{CellGroup}
\begin{CellGroup}
\Exercise{Problem 2}
\ExerciseText{Obtain the inverse DTFT of two lowpass filters

(a) \inlineTFinmath{{G_1}(\theta )\multsp =\multsp H\big(\frac{\theta }{2}\big)}.

(b) \inlineTFinmath{{G_2}(\theta )\multsp =\multsp H(2\theta )}.

Compare all three results. How does the width of the filter affect the time-domain characteristics of the unit sample response of the filter?}



\end{CellGroup}
\begin{CellGroup}
\Exercise{Problem 3}
\ExerciseText{(a) Plot the spectrum for \inlineTFinmath{M=10} and \inlineTFinmath{M=21}.

(b) Discuss how the spectrum of the finite, discrete-time pulse changes with increasing pulse width? Specifically, consider the mainlobe width and
sidelobe ripple.}



\end{CellGroup}

\end{CellGroup}
\begin{CellGroup}
\Section*{HOMEWORK PROBLEMS}
\begin{CellGroup}
\Subsubsection*{Homework Problem 15.1}
Given the following function of frequency (fundamental interval), calculate the inverse DTFT.

\dispSFinmath{\MathBegin{MathArray}{l}
L:=10  \\
\noalign{\vspace{1.25ex}}  \\
Y[\Mvariable{\theta \_}]:=\frac{\sin \big[\frac{\theta \multsp L}{2}\big]\multsp \exp \big[-\imag \frac{\theta }{2}\multsp (L-1)\big]}{\sin \big[\frac{\theta
}{2}\big]}
\MathEnd{MathArray}}
\begin{CellGroup}
\dispSFinmath{\MathBegin{MathArray}{l}
\Muserfunction{DisplayTogetherArray}[\{\Mfunction{Plot}[\Mfunction{Abs}[Y[\theta ]],\{\theta ,0,\pi \},  \\
\noalign{\vspace{0.5ex}}
\hspace{3.em} \Mvariable{PlotRange}->\Mvariable{All},\Mvariable{DisplayFunction}->\Mvariable{Identity}],\multsp \Mfunction{Plot}[\Mfunction{Arg}[Y[\theta
]],  \\
\noalign{\vspace{0.5ex}}
\hspace{3.em} \{\theta ,-\pi ,\pi \},\Mvariable{PlotRange}->\Mvariable{All},\Mvariable{DisplayFunction}->\Mvariable{Identity}]\}];\\
\MathEnd{MathArray}}
\mathGraphic{lecture15.tex_gr7.eps}

\end{CellGroup}

\end{CellGroup}
\begin{CellGroup}
\Subsubsection*{Homework Problem 15.2}
The summation defining the DTFT is infinite in length. In this problem explore the possibility of approximating the DTFT with a finite sum. Compute
such finite length "transforms" of the infinite-length sequence \inlineTFinmath{x[n]\multsp =\multsp {{\Big(\frac{1}{2}\Big)}^n}\multsp u[n]}. Try
a few different lengths of the summation. When does the graph of the finite length result become indistinguishable from the exact result?




\end{CellGroup}
\begin{CellGroup}
\Subsubsection*{Homework Problem 15.3}
Calculate and plot the DTFT of the three interesting signals defined below. Collectively, these are known as "windows" and they find application
in many discrete-time signal processing operations. Compare sidelobe ripple and mainlobe width of the following window sequences (L\(=\)2 M\(+\)1,
with M\(=\)9).

\dispSFinmath{\Muserfunction{win01}[\Mvariable{n\_},\Mvariable{M\_}]:=\Mfunction{If}\big[\multsp \Mfunction{Abs}[n]\leq M,\Big(1-\frac{\Mfunction{Abs}[n]}{M+1}\multsp
\Big)\multsp ,0\big]\multsp }
\dispSFinmath{\MathBegin{MathArray}{l}
\Muserfunction{win02}[\Mvariable{n\_},\Mvariable{M\_}]:=  \\
\noalign{\vspace{0.9375ex}}
\hspace{1.em} \Mfunction{If}\big[\multsp \Mfunction{Abs}[n]\leq M,\Big(\frac{7938}{18608}+\frac{9240\multsp }{18608}\cos \big[\frac{\multsp \pi \multsp
n}{M}\big]+\frac{1430\multsp }{18608}\cos \big[\frac{2\multsp \pi \multsp n}{M}\big]\Big)\multsp ,0\big]\\
\MathEnd{MathArray}}
\dispSFinmath{\Muserfunction{win03}[\Mvariable{n\_},\Mvariable{M\_}]:=\Mfunction{If}\big[\multsp \Mfunction{Abs}[n]\leq M,\frac{1}{2}\multsp \Big(1+\cos
\big[\frac{2\multsp \pi \multsp n}{2M+1}\big]\Big)\multsp ,0\big]}
\begin{CellGroup}
\dispSFinmath{\MathBegin{MathArray}{l}
\Muserfunction{DisplayTogetherArray}[\{\Mvariable{LinePlot}[  \\
\noalign{\vspace{0.5ex}}
\hspace{3.em} \Mfunction{Table}[\{n,\Muserfunction{win01}[n,9]\},\{n,-12,12\}],\Mvariable{PlotLabel}\rightarrow \Mvariable{Bartlett\multsp window}],
 \\
\noalign{\vspace{0.5ex}}
\hspace{2.em} \Muserfunction{LinePlot}[\Mfunction{Table}[\{n,\Muserfunction{win02}[n,9]\},\{n,-12,12\}],  \\
\noalign{\vspace{0.5ex}}
\hspace{3.em} \Mvariable{PlotRange}->\Mvariable{All},\Mvariable{PlotLabel}\rightarrow \Mvariable{Blackman\multsp window}],\Mvariable{LinePlot}[ 
\\
\noalign{\vspace{0.5ex}}
\hspace{3.em} \Mfunction{Table}[\{n,\Muserfunction{win03}[n,9]\},\{n,-12,12\}],\Mvariable{PlotLabel}\rightarrow \Mvariable{Hann\multsp window}]\}];\\
\MathEnd{MathArray}}
\mathGraphic{lecture15.tex_gr8.eps}

\end{CellGroup}

\end{CellGroup}
\begin{CellGroup}
\Subsubsection*{Homework Problem 15.4}
Consider an LTI system with impulse response

\dispSFinmath{h[\Mvariable{n\_}]:={{\Big(\frac{1}{2}\Big)}^{\Mfunction{Abs}[n]}}}
(a) Obtain the frequency response.

(b) Obtain the output \inlineTFinmath{y[n]} given the following input:



\dispSFinmath{x[\Mvariable{n\_}]:=\sin \big[\frac{3\multsp \pi \multsp n}{4}\big]}

\end{CellGroup}

\end{CellGroup}

\end{document}
